\section{Switching and finite holonomy}\label{sec:holonomy}

Write $a_i$ for the sign of $\{i,i+1\}$ and $d_i$ for the sign of
$\{i,i+2\}$, with cyclic indices. The triangle sign and Hamilton-cycle
holonomy are
\[
 \tau_i=a_i a_{i+1}d_i,\qquad \alpha=\prod_{i=0}^{n-1}a_i.
\]
Switching by vertex signs $g_i\in\{\pm1\}$ replaces $A$ by $GAG$,
where $G=\diag(g_i)$. It preserves the spectrum, every $\tau_i$, and
$\alpha$.

\begin{lemma}\label{lem:gauge}
For prescribed triangle signs $(\tau_i)_{i\in\Z/n\Z}$, there are exactly
two switching classes, indexed by $\alpha\in\{\pm1\}$. Each has a unique
representative with
\[
 a_0=\cdots=a_{n-2}=1,\qquad a_{n-1}=\alpha,
 \qquad d_i=\tau_i a_i a_{i+1}.
\]
\end{lemma}
\begin{proof}
Set $g_0=1$ and recursively $g_{i+1}=g_i a_i$ for $0\le i<n-1$.
After switching, the first $n-1$ step-one signs are positive. Their product
with the remaining sign is $\alpha$, so that remaining sign is $\alpha$.
The triangle equations force every step-two sign. Conversely, the displayed
signs realize the prescribed data for either choice of $\alpha$.
A switch between two such representatives is constant along the path
$0,1,\ldots,n-1$, hence constant on all vertices and acts trivially.
Finally, distinct values of the switching invariant $\alpha$ cannot be
switching equivalent.
\end{proof}

For the repeated word $t^m$ on $N=8m$ vertices, write
$A_m^\pm=A_{8m}(t^m,\pm1)$. The positive representative has all step-one
signs positive and step-two signs $t_{i\bmod8}$. The negative representative
reverses exactly the three distinct edges
\begin{equation}
 \{N-1,0\},\qquad\{N-2,0\},\qquad\{N-1,1\}.
 \label{eq:seam}
\end{equation} Reversing only the Hamilton edge would
change two triangle signs. The two additional step-two reversals are
therefore essential.

Extend $u\in\C^{8m}$ to the integer line by
\begin{equation}
 u_{i+8m}=\alpha u_i.
 \label{eq:extension}
\end{equation}
Let $t_i=t_{i\bmod8}$. In either seam gauge the adjacency action is the
restriction of
\begin{equation}
 (Lu)_i=u_{i-1}+u_{i+1}+t_{i-2}u_{i-2}+t_i u_{i+2}.
 \label{eq:infinite-action}
\end{equation}
In particular, the signs of all cut-crossing edges in the finite matrix
agree with \eqref{eq:extension}, including for $m=1$.

\begin{lemma}\label{lem:bloch}
The complexification of $A_m^\alpha$ is unitarily equivalent to
\[
 \bigoplus_{z^m=\alpha}H(z),
\]
where, in the order $0,1,\ldots,7$,
\begin{equation}
 H(z)=\begin{pmatrix}
0&1&1&0&0&0&z^{-1}&z^{-1}\\
1&0&1&1&0&0&0&-z^{-1}\\
1&1&0&1&-1&0&0&0\\
0&1&1&0&1&1&0&0\\
0&0&-1&1&0&1&-1&0\\
0&0&0&1&1&0&1&-1\\
z&0&0&0&-1&1&0&1\\
z&-z&0&0&0&-1&1&0
\end{pmatrix}.
\label{eq:fiber}
\end{equation}
\end{lemma}
\begin{proof}
For each root $z^m=\alpha$ and $v\in\C^8$, set
$u_{8j+a}=m^{-1/2}z^jv_a$ for $0\le j<m$ and $0\le a<8$.
This satisfies \eqref{eq:extension}. Substitution into
\eqref{eq:infinite-action} gives \eqref{eq:fiber}. For two different
allowed phases $z,w$, the ratio $z\overline w$ is a nontrivial $m$th
root of unity, so
\[
 \sum_{j=0}^{m-1}(z\overline w)^j=0.
\]
Thus the $m$ fiber spaces are mutually orthogonal, each has dimension
eight, and together they exhaust dimension $8m$. The displayed map is
unitary. Since $|z|=1$, the fiber matrix is Hermitian.
\end{proof}
